% SEGMENT OLP-0442-S01
% Part: normal-modal-logic
% Chapter: completeness
% Section: introduction

\documentclass[../../../include/open-logic-section]{subfiles}

\begin{document}

\olfileid{nml}{com}{int}

\olsection{அறிமுகம்}

$\Sigma$ ஒரு வாய்ப்புநிலை அமைப்பு எனில்,~$\Sigma$ விலுள்ள எல்லா
!!{formula}s இன் எல்லா நிகழ்வுகளும் செல்லுபடியாகும் மாதிரிகளின் எந்த
வகுப்பு~$\mClass{C}$ இலும், $\Sigma \Proves !A$ எனில்~$!A$
செல்லுபடியாகும் என்பதைச் சரித்தன்மைத் தேற்றம் நிறுவுகிறது. குறிப்பாக,
$\Log{K} \Proves !A$ எனில் எல்லா மாதிரிகளிலும்~$!A$ மெய்யாகும்;
$\Log{KT} \Proves !A$ எனில் எல்லாத் தற்சுட்டு மாதிரிகளிலும்~$!A$
மெய்யாகும்; $\Log{KD} \Proves !A$ எனில் எல்லாத் தொடர் மாதிரிகளிலும்
$!A$ மெய்யாகும்; இவ்வாறே தொடரும்.

நிறைவு என்பது சரித்தன்மையின் மறுதலையாகும்:~\Log{K} நிறைவானது என்பது,
எடுத்துக்காட்டாக, !!a{formula}~$!A$ செல்லுபடியாக இருந்தால் $\Proves !A$
என்பதாகும். சரித்தன்மையை நிறுவுவதைவிட நிறைவை நிறுவுவது மிகவும் கடினம்.
முதலில் எதிர்நிலைமாற்றத்தைக் கருதுவது பயனுள்ளது: $\Proves/ !A$ ஆகும்
போதெல்லாம் ஓர் எதிர்மாதிரி, அதாவது $\mSat/{M}{!A}$ ஆகும்
மாதிரி~$\mModel{M}$ ஒன்று இருந்தால், இருந்தால் மட்டுமே,~\Log{K}
நிறைவானது. நிகர்வகையில் ($!A$ ஐ மறுத்து), $\Proves/ \lnot !A$ ஆகும்
போதெல்லாம்~$!A$ க்கு ஒரு மாதிரி உண்டு என்று நிறுவலாம். அத்தகைய மாதிரியைக்
கட்டமைக்கும்போது, $!A$ இல் அடங்கியுள்ள தகவலைப் பயன்படுத்தலாம். குறிப்பிட்ட
!!{formula}s க்கான மாதிரிகளைக் காணும்போதும் இதையே அடிக்கடி செய்கிறோம்:
எடுத்துக்காட்டாக, $p \lif \Box q$ க்கு ஓர் எதிர்மாதிரியைக் காண விரும்பினால்,
$p$ மெய்யாகவும் $\Box q$ பொய்யாகவும் இருக்கும் ஓர் உலகை அது கொண்டிருக்க
வேண்டும் என்று அறிவோம். $\Box q$ பொய்யாகும் ஓர் உலகம் என்றால், அதிலிருந்து
அணுகத்தக்கதும்~$q$ பொய்யாக இருப்பதுமான ஓர் உலகம் இருக்க வேண்டும்.
நாம் அறிய வேண்டியது அவ்வளவே: எந்த உலகுகள் எந்த !!{propositional variable}s
ஐ மெய்யாக்குகின்றன, எந்த உலகுகளிலிருந்து எந்த உலகுகள் அணுகத்தக்கவை.

ஆனால் நிறைவை நிறுவும் நேர்வில், நாம் மாதிரி கட்டமைக்கும் ஒரு குறிப்பிட்ட
!!{formula}~$!A$ இல்லை. $\Proves/[\Sigma] \lnot !A$ ஆக இருக்கும் ஒவ்வொரு
$!A$ க்கும் ஒரு மாதிரி உண்டு என்று நிறுவ விரும்புகிறோம். இது ஒரு
குறைந்தபட்சக் கோரிக்கை; ஏனெனில் \emph{$\Proves[\Sigma] \lnot !A$ எனில்},
சரித்தன்மையின்படி ($\Sigma$ மெய்யாக இருக்கும்)~$!A$ கான மாதிரி எதுவும்
இல்லை. இப்போது $\Proves/[\Sigma] \lnot !A$ என்பது, அப்போதும் அப்போது
மட்டுமே, $!A$ ஆனது~$\Sigma$-முரணற்றது என்பதைக் கவனிக்க.
% SOURCE CORRECTION TA-NML-025: the equivalent non-derivability statement has no premise set; Sigma already indexes the modal system.
($\Proves/[\Sigma] \lnot !A$, $!A \Proves/[\Sigma] \lfalse$ ஆகியவை
நிகரானவை என்பதை நினைவுகூர்க.) ஆகவே ஒவ்வொரு $\Sigma$-முரணற்ற
!!a{formula} க்கும் ஒரு மாதிரியைக் கட்டமைப்பதே நமது பணி.

நாம் பயன்படுத்தவிருக்கும் தந்திரம்,~$!A$ ஐக் கொண்டிருப்பதோடு,~$!A$ ஐ
மெய்யாக்கும் உலகம் எவ்வாறு இருக்க வேண்டும் என்று நமக்குக் கூறும் பிற
வாய்பாடுகளையும் கொண்ட ஒரு $\Sigma$-முரணற்ற !!{formula}s கணத்தைக்
கண்டுபிடிப்பதாகும். அத்தகைய கணங்கள் \emph{நிறைவான} $\Sigma$-முரணற்ற
கணங்கள். $!A$ ஐ மெய்யாக்க ஒற்றை உலகைக் கொண்ட ஒரு மாதிரியைக் கட்டமைப்பது
போதாது; அது பல உலகுகளையும் ஓர் அணுகல்தன்மைத் தொடர்பையும் கொண்டிருக்க
வேண்டும். $!A$ ஐக் கொண்ட நிறைவான $\Sigma$-முரணற்ற கணம், $\Box !B$,
$\Diamond !C$ ஆகிய வடிவங்களிலுள்ள பிற !!{formula}s ஐயும் கொண்டிருக்கும்.
அணுகத்தக்க எல்லா உலகுகளிலும்~$!B$ மெய்யாக இருக்க வேண்டும்; குறைந்தது
ஒன்றிலாவது~$!C$ மெய்யாக இருக்க வேண்டும். இதைச் செய்ய, சாத்தியமான
\emph{எல்லா} நிறைவான $\Sigma$-முரணற்ற கணங்களையும் உலகுகளின் கணத்திற்கான
அடிப்படையாக எடுத்துக்கொள்வோம். நமது மாதிரியில் ஒரு நிறைவான
$\Sigma$-முரணற்ற கணம் மற்றொன்றிலிருந்து எப்போது அணுகத்தக்கதாக எண்ணப்பட
வேண்டும் என்பதைக் கண்டறிவது ஒரு நுட்பமான பகுதியாக இருக்கும்.

இவ்வாறு வரையறுக்கப்பட்ட மாதிரியில், ஒரு உலகில்---அதுவும் ஒரு நிறைவான
$\Sigma$-முரணற்ற கணமே---$!A$ மெய்யாகும் என்பது, அப்போதும் அப்போது மட்டுமே,
$!A$ அந்தக் கணத்தின் !!a{element} ஆகும் என்று காட்டுவோம். $!A$ ஆனது
$\Sigma$-முரணற்றது எனில், அது குறைந்தது ஒரு நிறைவான $\Sigma$-முரணற்ற
கணத்தின் !!a{element} ஆக இருக்கும் (இவ்வுண்மையை நிறுவுவோம்); ஆகவே~$!A$
மெய்யாகும் ஓர் உலகம் இருக்கும். எனவே, ஒவ்வொரு $\Sigma$-முரணற்ற
!!{formula}~$!A$ உம் ஏதோவோர் உலகில் மெய்யாக இருக்கும் ஒரேயொரு மாதிரி
நமக்குக் கிடைக்கும். இந்த ஒரே மாதிரியே~$\Sigma$ க்கான \emph{நியம}
மாதிரியாகும்.

\end{document}
